% !TeX root = ../../Book3.tex
\section{模函数与 \texorpdfstring{\(j\)}{j}-不变量}
\label{b3:sec:2306}

\(E_4\) 与 \(E_6\) 随格基的改变而乘上不同因子，不能直接用它们的数值比较两个复环面。若先把权调成相同，再取商，因子便会约去。问题随之变为：能否找到一个这样的商，使它的值不仅在等价的格上相同，而且足以辨认全部格比例类？下面将用价公式回答这个逆向问题，再给所得参数空间建立具体的复坐标。

\subsection{模函数}
\label{b3:subsec:230601}

\begin{definition}\label{b3:def:23-modular-function}
设 \(\mathfrak H=\{\tau\in\mathbb C:\operatorname{Im}\tau>0\}\)，群 \(\mathrm{SL}_2(\mathbb Z)\) 按分式线性变换作用，\(h\) 为 \(\mathfrak H\) 上的亚纯函数。若
\[
 h\left(\frac{a\tau+b}{c\tau+d}\right)=h(\tau)
 \quad\left(\begin{pmatrix}a&b\\c&d\end{pmatrix}\in\mathrm{SL}_2(\mathbb Z)\right),
\]
且周期函数 \(h\) 的 \(q=e^{2\pi\mathrm i\tau}\) 表达式在 \(q=0\) 亚纯，则称 \(h\) 为全模群的\term[mo-hanshu]{模函数}[modular function]。这里的亚纯函数不包括恒取 \(\infty\) 的映射，允许恒零函数。
\end{definition}

尖点条件允许有限个负次幂；它不要求在尖点有界。若 \(h\not\equiv0\)，则在充分高的半平面有
\[
 h(\tau)=q^m\bigl(a_m+a_{m+1}q+\cdots\bigr),
 \qquad m\in\mathbb Z,\quad a_m\ne0.
\]
正的 \(m\) 表示尖点零点，负的 \(m\) 表示尖点极点。由于 \(|q|=e^{-2\pi\operatorname{Im}\tau}\)，一个 \(r\) 阶尖点极点沿高半平面的增长量级为 \(e^{2\pi r\operatorname{Im}\tau}\)。因此“在无穷远亚纯”说的是指数坐标中的有限阶行为，不能照搬平面坐标 \(1/\tau\) 下的极点定义。

\ProseParagraphBreak
若 \(f,g\in M_k\) 且 \(g\not\equiv0\)，则 \(f/g\) 为模函数。上半平面中的变换因子相消；在任一有限点把分子、分母各自提出其零点幂次，商只有有限阶零点或极点。尖点处同样对两个收敛的 \(q\) 幂级数作局部分解。这说明商在有限点和尖点都亚纯，不能仅凭形式上“权相减为零”省略这两处检查。

\ProseParagraphBreak
非零模函数的倒数、两个模函数的和与积仍为模函数，因而它们组成域。本章稍后将把这个域确定为 \(\mathbb C(j)\)。这里 \(\mathbb C(j)\) 指常系数有理函数代入 \(j\) 所得的函数域，而 \(M_*\) 是按权分次的全纯模形式环；前者允许零权函数有极点，后者用非零权容纳丰富的全纯对象，两者承担不同职责。

\subsection{\texorpdfstring{\(j\)}{j}-不变量}
\label{b3:subsec:230602}

判别式 \(\Delta\) 在上半平面没有零点，且在尖点恰好一阶消失。因此用它作分母，可以把全部极点集中到唯一的尖点。分子须有权十二；选 \(E_4^3\) 还使首项系数为一，这固定了一个方便的归一化。

\begin{definition}\label{b3:def:23-j-invariant}
对 \(\tau\in\mathfrak H\)，令 \(E_4,E_6\) 为全模群上常数项归一为一的权四、权六 Eisenstein 级数，并令 \(\Delta=(E_4^3-E_6^2)/1728\)。定义
\term[j-bubianliang]{\(j\)-不变量}[j-invariant]为
\[
 j(\tau)=\frac{E_4(\tau)^3}{\Delta(\tau)}.
\]
\end{definition}
\symidx{\(j(\tau)\)：椭圆模不变量}

两个权十二因子相除，故 \(j\) 对全模群不变。由前节的 Fourier 展开，
\[
 E_4^3=1+720q+179280q^2+O(q^3),\qquad
 \Delta=q\bigl(1-24q+252q^2+O(q^3)\bigr).
\]
先求第二个括号的倒数：若写成 \(1+b_1q+b_2q^2+O(q^3)\)，乘回原式比较系数得到 \(b_1=24\)、\(b_2=24^2-252=324\)。再与分子相乘，得到
\begin{equation}\label{b3:eq:23-j-fourier}
 j=q^{-1}+744+196884q+O(q^2).
\end{equation}
这里的常数项为 \(720+24\)，一次项为 \(179280+720\cdot24+324\)。归一化中的 \(1728\) 使尖点主部恰为 \(q^{-1}\)；若改取 \(j/1728\)，分类性质不变，但特殊值和全部展开系数都须相应改变。

\Needspace{6\baselineskip}
\begin{theorem}[\texorpdfstring{\(j\)}{j} 对模群轨道的分离]\label{b3:thm:23-j-orbits}
设 \(j=E_4^3/\Delta\)，其中 \(E_4,E_6\) 为全模群归一化 Eisenstein 级数，\(\Delta=(E_4^3-E_6^2)/1728\)，并令 \(\rho=e^{2\pi\mathrm i/3}\)。对每个 \(w\in\mathbb C\)，集合
\(\{\tau\in\mathfrak H:j(\tau)=w\}\) 恰为一条 \(\mathrm{PSL}_2(\mathbb Z)\) 轨道。若 \(p\) 在这条轨道上，\(j-w\) 在上半平面参数 \(\tau-p\) 中的零点阶数为
\[
 \begin{cases}
 3,&p\text{ 与 }\rho\text{ 模群等价},\\
 2,&p\text{ 与 }\mathrm i\text{ 模群等价},\\
 1,&p\text{ 为其他点}.
 \end{cases}
\]
此外 \(j(\rho)=0\)、\(j(\mathrm i)=1728\)。
\end{theorem}
\begin{proof}
固定 \(w\)，考察
\[
 F_w=E_4^3-w\Delta\in M_{12}.
\]
它的常数项为一，故非恒零，且 \(v_\infty(F_w)=0\)。由于 \(\Delta\) 内部无零，\(F_w\) 与 \(j-w\) 在上半平面具有相同零点及阶数。\cref{b3:thm:23-valence}于是给出
\[
 \frac12v_{\mathrm i}(F_w)+\frac13v_\rho(F_w)
 +\sum_{p\ne\mathrm i,\rho}v_p(F_w)=1,
\]
普通点仍按轨道取一个代表。

这个等式本身只给加权总数，必须再知道每条非空零点轨道至少贡献一，才能推出唯一性。普通点的零点阶数为正整数，已有这一性质。在椭圆固定点使用 \(j-w\) 的不变性：由\cref{b3:prop:23-elliptic-orders}的权零情形，它在 \(\mathrm i\) 的阶数为二的倍数，在 \(\rho\) 的阶数为三的倍数。因此这两个位置除以稳定子阶数后的贡献也为正整数。

总数为一首先排除了零点集合为空，又只允许一条零点轨道，并强迫其商参数中的阶数为一。换回上半平面参数，就得到所述三种阶数。最后，由\cref{b3:cor:23-eisenstein-zeros}，\(E_4(\rho)=0\)，而 \(E_6(\mathrm i)=0\)。代入
\[
 j=\frac{E_4^3}{\Delta},\qquad
 j-1728=\frac{E_6^2}{\Delta}
\]
便得到两个特殊值。
\end{proof}

\ProseParagraphBreak
此证明同时给出存在和唯一性。不能先假设某个 \(w\) 被取到，再只证明原像唯一；价公式右侧的正数保证每个有限 \(w\) 都确实有原像。它也说明基本域边界上的两个代表不应算成两个复结构：总数始终按群轨道计，而非按闭基本域中互异点的个数计。

\ProseParagraphBreak
例如 \(E_4^3-1728\Delta=E_6^2\) 的零点全部位于 \(\mathrm i\) 的轨道，每处为二阶；\(E_4^3\) 的零点位于 \(\rho\) 的轨道，每处为三阶。若 \(w\notin\{0,1728\}\)，则 \(E_4^3-w\Delta\) 只有一条普通零点轨道，每处为简单零点。三个情形的加权总数都是一。

\subsection{复环面的 \texorpdfstring{\(j\)}{j} 分类}
\label{b3:subsec:230603}

前面的轨道结论是在参数 \(\tau\) 上得到的。要把它解释成复环面的分类，还须排除一种可能：两个格不成比例，却有某种非线性的全纯坐标变化把它们的商曲面同构起来。平面覆盖使这一问题可以转为整函数的导数估计。

\Needspace{6\baselineskip}
\begin{proposition}[复环面之间的全纯映射]\label{b3:prop:23-torus-affine}
设 \(\Lambda,\Lambda'\subset\mathbb C\) 为秩二格，\(T_\Lambda=\mathbb C/\Lambda\)、\(T_{\Lambda'}=\mathbb C/\Lambda'\) 带有平面商坐标所定义的复结构。每个全纯映射 \(F:T_\Lambda\to T_{\Lambda'}\) 都可写为
\[
 F([z]_\Lambda)=[az+b]_{\Lambda'},\qquad
 a,b\in\mathbb C,\quad a\Lambda\subseteq\Lambda'.
\]
反之，每个满足上述包含条件的仿射式都定义全纯映射。映射 \(F\) 双全纯当且仅当 \(a\ne0\) 且 \(a\Lambda=\Lambda'\)；若还要求 \(F([0])=[0]\)，则可取 \(b=0\)。
\end{proposition}
\begin{proof}
记商投影为 \(\pi_\Lambda,\pi_{\Lambda'}\)。\cref{b3:prop:22-torus-charts}已证明这两个投影是全纯覆盖：取足够小的圆盘，其不同格平移互不相交，投影在每一片上双全纯。平面单连通，故\cref{b3:thm:11-map-lifting}应用于 \(F\circ\pi_\Lambda:\mathbb C\to T_{\Lambda'}\)，给出连续提升 \(\widetilde F:\mathbb C\to\mathbb C\)。提升在局部等于 \(\pi_{\Lambda'}\) 的全纯逆分支与 \(F\circ\pi_\Lambda\) 的复合，因此为整函数。

对每个 \(\lambda\in\Lambda\)，两个点 \(z\) 与 \(z+\lambda\) 的投影相同，所以
\[
 \widetilde F(z+\lambda)-\widetilde F(z)\in\Lambda'.
\]
左端随 \(z\) 连续，右端是离散集合；平面连通使这个差对 \(z\) 为常数。对 \(z\) 求导，便知 \(\widetilde F'\) 为整的 \(\Lambda\)-周期函数。它在闭基本平行四边形上有界，再由周期性在整个平面有界，Liouville 定理使 \(\widetilde F'=a\) 为常数。积分得到 \(\widetilde F(z)=az+b\)，原来的差式给 \(a\lambda\in\Lambda'\)。

反向的良定义也正由这个条件保证：改变代表 \(z\mapsto z+\lambda\)，像代表只增加 \(a\lambda\in\Lambda'\)；局部商坐标中是仿射全纯函数。

若 \(F\) 双全纯，便不恒定，所以 \(a\ne0\)。若有 \(\lambda'\in\Lambda'\setminus a\Lambda\)，则 \([\lambda'/a]_\Lambda\ne[0]_\Lambda\)，却与 \([0]_\Lambda\) 映到同一点，违背单射。因此 \(a\Lambda=\Lambda'\)。反之，等式成立时 \(w\mapsto(w-b)/a\) 良定义地下降为逆映射，故双全纯。若 \(F([0])=[0]\)，则 \(b\in\Lambda'\)，在商中可以去掉这一项。
\end{proof}

这里并未预先假设 \(F\) 保持加法。结论说明，先用目标环面的平移把 \(F([0])\) 移回原点，所得全纯映射自动由复线性映射下降，因此自动是群同态。若只要求实线性同构，则任意两组实基都能相互映到，全部秩二格的商在拓扑上都是环面；复线性这一条件才保留共形结构。

\Needspace{6\baselineskip}
\begin{theorem}[复环面的模参数分类]\label{b3:thm:23-torus-j-classification}
设 \(\Lambda=\omega_1(\mathbb Z+\mathbb Z\tau)\)、\(\Lambda'=\omega_1'(\mathbb Z+\mathbb Z\tau')\)，其中 \(\omega_1,\omega_1'\ne0\)，\(\tau,\tau'\in\mathfrak H\)。以下条件等价：
\begin{enumerate}
\item 复环面 \(\mathbb C/\Lambda\) 与 \(\mathbb C/\Lambda'\) 双全纯；
\item 存在 \(a\in\mathbb C^*\) 使 \(a\Lambda=\Lambda'\)；
\item \(\tau'\) 与 \(\tau\) 属于同一 \(\mathrm{PSL}_2(\mathbb Z)\) 轨道；
\item \(j(\tau')=j(\tau)\)。
\end{enumerate}
每个有限复数均为某个复环面的 \(j\) 值。
\end{theorem}
\begin{proof}
前两项等价由\cref{b3:prop:23-torus-affine}给出，其中常数平移不改变格条件。\cref{b3:prop:23-lattice-orbits}已用正向整数基变换证明：归一化格 \(\mathbb Z+\mathbb Z\tau\) 与 \(\mathbb Z+\mathbb Z\tau'\) 成复比例，当且仅当两参数由行列式一的整数矩阵联系；缩放因子 \(\omega_1,\omega_1'\) 不影响这一判别。最后两项等价以及有限值的满射性由\cref{b3:thm:23-j-orbits}得到。
\end{proof}

\ProseParagraphBreak
对格 \(\Lambda_\tau\)，\secref{b3:sec:2304}给出的常数换算使
\[
 j(\tau)=1728\frac{g_2(\Lambda_\tau)^3}
 {g_2(\Lambda_\tau)^3-27g_3(\Lambda_\tau)^2}.
\]
分子、分母在格缩放下同乘十二次逆幂，故也可直接用任意未归一化的格计算这个比值。方格的 \(g_3=0\) 给 \(j=1728\)，六角格的 \(g_2=0\) 给 \(j=0\)。特殊值表示额外的旋转对称，但两类环面仍各自只有一个同构类，分类定理没有在这些位置失效。

\ProseParagraphBreak
若 \(Y,Y'>1\) 且 \(Y\ne Y'\)，则 \(\mathrm iY,\mathrm iY'\) 是基本域内部不同点，因此对应的矩形环面不双全纯，\(j(\mathrm iY)\ne j(\mathrm iY')\)。另一方面，\(\mathrm iY\) 与 \(\mathrm i/Y\) 由 \(S\) 联系，代表同一环面。这个例子说明“参数不同”与“复结构不同”之间还隔着换基。

\ProseParagraphBreak
由\cref{b3:thm:22-uniformization}，每个格环面已经与它的 Weierstrass 光滑射影三次曲线双全纯对应，原点对应无穷远点。因此上述结论也分类这些由格构造的曲线的保原点解析同构。\(j\) 相等意味着相应复结构同构，并不要求两条仿射方程的系数逐项相等。

\subsection{模曲线 \texorpdfstring{\(X(1)\)}{X(1)} 的几何解释}
\label{b3:subsec:230604}

目前 \(j\) 已给出轨道集合与 \(\mathbb C\) 的双射。集合双射还没有说明连续性、局部复坐标和无穷远点。为理解“模参数空间也是一张 Riemann 曲面”，下面把这些结构逐一建立。

令 \(G=\mathrm{PSL}_2(\mathbb Z)\)，\(Y(1)=G\backslash\mathfrak H\) 带商拓扑，\(\pi:\mathfrak H\to Y(1)\) 为投影。\secref{b3:sec:2301}已证明紧集有限相交性质及局部商模型。普通点附近可取投影单射的小圆盘；在 \(p=\mathrm i\) 或 \(\rho\) 附近，先用
\[
 u_p(\tau)=\frac{\tau-p}{\tau-\overline p}
\]
把稳定子作用化为旋转，再用 \(u_p^e\) 作商坐标，分别取 \(e=2,3\)。其他与这两点等价的位置先经群元搬回即可。这是\cref{b3:prop:23-stabilizer-chart}的具体内容；不能在固定点直接把 \(\pi\) 当成无分歧覆盖。

\ProseParagraphBreak
这些图之间的转移是全纯的。在没有固定点的小片上，选定局部提升后，转移由 Möbius 变换及全纯坐标复合给出；椭圆点中心与普通小片不会同时相交，在去心重叠处可选 \(u_p^e\) 的一个根分支。同一椭圆轨道可固定一个上述坐标图；缩小此图或改用等价代表时，群变换与稳定子旋转相容，相应不变全纯函数按幂级数下降到商坐标。因此这些局部图给出相容的复结构。

\Needspace{6\baselineskip}
\begin{proposition}[唯一尖点的紧化]\label{b3:prop:23-modular-compactification}
设 \(G=\mathrm{PSL}_2(\mathbb Z)\) 按分式线性变换作用于 \(\mathfrak H\)。在 \(Y(1)=G\backslash\mathfrak H\) 中加入一个点 \(\infty\)，规定它的邻域基为
\[
 U_Y=\{\infty\}\cup\pi\{\tau\in\mathfrak H:\operatorname{Im}\tau>Y\},
 \qquad Y>1.
\]
用 \(q=e^{2\pi\mathrm i\tau}\) 将 \(U_Y\) 识别为 \(|q|<e^{-2\pi Y}\)，并以 \(q=0\) 对应新增点。它与普通点坐标、椭圆点坐标 \(u_p^e\) 共同定义一个紧连通 Riemann 曲面，记为 \(X(1)\)。
\end{proposition}
\begin{proof}
\cref{b3:prop:23-cusp-chart}保证当 \(Y>1\) 时，高半平面内只有整数平移可能识别两个点。因此 \(q\) 在 \(U_Y\setminus\{\infty\}\) 上良定义且为一一的局部全纯坐标，像恰为去心圆盘。补上零点就得到完整圆盘。在两个尖点图的重叠上坐标相同；与内部图的重叠上，\(q\) 的导数 \(2\pi\mathrm i q\) 非零，局部对数给出全纯逆，故转移相容。

先证明内部轨道商是 Hausdorff 的。取不属于同一轨道的 \(p,q\in\mathfrak H\)，并取它们的紧邻域 \(K,L\)。由\cref{b3:prop:23-proper-action}，满足 \(\gamma K\cap L\ne\varnothing\) 的群元只有有限个。每个这样的群元均有 \(\gamma p\ne q\)，故可缩小 \(p,q\) 的开邻域，使其在这个群元下不相交；有限次取交后，得到 \(U\subset K,V\subset L\)，使所有这些群元均满足 \(\gamma U\cap V=\varnothing\)。候选集合以外的群元本来就不可能相交。因此 \(\pi(U),\pi(V)\) 是不交的开邻域：商投影为开映射，因为一个开集全部群平移的并仍开。

还须把新尖点与任一内部点分开。固定 \(\tau_0\in\mathfrak H\)，选一个含它的小紧邻域 \(K\)，其中 \(\eta\le\operatorname{Im}\tau\le B\)、\(\eta>0\)。对 \(\tau=x+\mathrm i y\in K\) 及任意群元，若 \(c\ne0\)，则
\[
 \operatorname{Im}(\gamma\tau)
 =\frac{y}{(cx+d)^2+c^2y^2}\le\frac1{c^2y}\le\eta^{-1};
\]
若 \(c=0\)，则虚部仍为 \(y\le B\)。取 \(Y>\max\{1,B,\eta^{-1}\}\)，\(\pi(K)\) 便与 \(U_Y\) 不交，从而尖点与内部点也有不交邻域。

将上半平面的可数基用开商投影送到 \(Y(1)\)，可得 \(Y(1)\) 的可数基。再添上整数高度的 \(U_Y\)，得到 \(X(1)\) 的可数基。因此上述局部图给出 Hausdorff、第二可数的一维复流形。\(Y(1)\) 是连通上半平面的连续像，而每个尖点邻域都与它相交，故添点后仍连通。

最后取 \(Y>1\)。截断基本域
\[
 K_Y=\{\tau\in\mathfrak H:|\operatorname{Re}\tau|\le1/2,
 \ |\tau|\ge1,\ \operatorname{Im}\tau\le Y\}
\]
是紧集：它闭有界，且虚部至少为 \(\sqrt3/2\)。它的投影覆盖高度不超过 \(Y\) 的基本域部分。其余部分加上尖点，在一个稍大尖点图中恰由闭圆盘 \(|q|\le e^{-2\pi Y}\) 参数化，也是紧的。这两个紧集的并覆盖 \(X(1)\)，故 \(X(1)\) 紧致。
\end{proof}

这个紧曲面称为全模群的\term[mo-quxian]{模曲线}[modular curve] \(X(1)\)。此处只讨论这一特定群的解析商及其紧化，不需要一般模空间的代数几何定义。单独加入的尖点不对应一个秩二格环面；当归一化基本域中的虚部趋向无穷时，它记录的是离开所有紧参数集的极限端。

\ProseParagraphBreak
模函数恰好下降为 \(X(1)\) 上的亚纯函数。普通点的下降来自群不变性；在阶数为 \(e\) 的椭圆点，局部 Laurent 级数在 \(u\mapsto\zeta_eu\) 下不变，故只含 \(u^{en}\)，改写成 \(u^e\) 的 Laurent 级数后仍只有有限个负幂。尖点处的亚纯性正是定义中的 \(q\) 条件。反过来，曲面上的亚纯函数拉回 \(\mathfrak H\) 后不变，并在尖点给出同样的有限负幂展开。这一步还说明：有限点在上半平面中的阶数，要除以稳定子阶数，才是商曲面上的阶数。

\Needspace{6\baselineskip}
\begin{theorem}[模曲线的球面坐标]\label{b3:thm:23-j-sphere}
令 \(X(1)\) 为全模群的上述紧化商曲面，\(j=E_4^3/\Delta\)，其中 \(E_4,E_6\) 为全模群归一化 Eisenstein 级数，\(\Delta=(E_4^3-E_6^2)/1728\)。令 \(j(\infty)=\infty\)，则 \(j\) 诱导双全纯映射
\[
 j:X(1)\longrightarrow\widehat{\mathbb C\kern0pt}.
\]
因此全模群的模函数域为 \(\mathbb C(j)\)。其中在整个上半平面全纯的模函数恰为 \(j\) 的多项式，若在尖点也全纯，则必为常数。
\end{theorem}
\begin{proof}
在内部，\cref{b3:thm:23-j-orbits}证明每个有限值恰对应一个商点，并在商坐标中为简单取值。换言之，\(j-j(p)\) 在普通坐标或 \(u_p^e\) 中都有非零一次系数，局部逆函数定理给全纯逆。尖点处，由\eqref{b3:eq:23-j-fourier}，
\[
 \frac1j=q\bigl(1-744q+O(q^2)\bigr),
\]
一次系数也非零；目标无穷远坐标是 \(1/j\)，故此处同样有全纯局部逆。内部没有极点，尖点是无穷远值的唯一原像。因此全局双射及各处全纯局部逆拼成所述双全纯映射。

若 \(h\) 是模函数，它下降到 \(X(1)\) 后与 \(j^{-1}\) 复合，是球面上的亚纯函数。\cref{b3:thm:03-sphere-rational}说明这种函数必为有理函数：球面紧致使极点只有有限多个，减去有限点的主要部分后，只剩一个在无穷远至多有极点的整函数，即多项式。因此 \(h=R(j)\)，其中 \(R\in\mathbb C(t)\)。反向把任意有理函数代入这个球面坐标，当然给出模函数。

若 \(h\) 在上半平面无极点，\(R\) 就不能有有限极点，因为每个有限值都被 \(j\) 取到；有理函数没有有限极点，当且仅当它是多项式。若它还在尖点无极点，该多项式在无穷远也全纯，只能为常数。
\end{proof}

\ProseParagraphBreak
例如一个模函数若只在尖点有至多二阶极点，就具有 \(Aj^2+Bj+C\) 的形式。\(j^2=q^{-2}+1488q^{-1}+\cdots\)，因此先用 \(j^2\) 匹配二阶主部，再用 \(j\) 匹配一阶主部，剩余仅是常数。这样的有限消元由函数域定理保证，不是仅凭若干 Fourier 系数作猜测。

\ProseParagraphBreak
尖点亚纯性在这里不可删除。函数 \(e^j\) 在上半平面全纯且不变，但在尖点有本性奇点：因为 \(1/j\) 是该点的一个全纯局部坐标，\(e^j\) 在此坐标中就是 \(e^{1/z}\)。它既不是模函数，也不是 \(j\) 的有理函数。另一方面，\(j\) 自己已经说明“内部全纯且不变”不强迫常值；必须再控制唯一缺失的尖点。

\ProseParagraphBreak
\(j\) 在上半平面有二阶、三阶局部分歧，与它在 \(X(1)\) 上双全纯并不冲突。先商掉稳定子之后，平方或立方参数已经吸收这些重数，\(j\) 在商坐标中只剩一阶变化。至此，格的换基、模形式的权、价公式中的分数权重和复环面的一个复数参数，便由明确的局部计算衔接起来。
